\documentclass{article}
\usepackage{amsmath}
\usepackage{amssymb}
\usepackage{amsthm}
\usepackage{mathtools}
\usepackage{geometry}
\usepackage{hyperref}

\newcommand{\rk}{\operatorname{rank}}
\newcommand{\GL}{\operatorname{GL}}
\newcommand{\R}{\mathbb{R}}
\newcommand{\Mat}{\mathrm{Mat}}
\newcommand{\rad}{\operatorname{rad}}

\theoremstyle{definition}
\newtheorem{theorem}{Theorem}[section]
\newtheorem{proposition}[theorem]{Proposition}
\newtheorem{lemma}[theorem]{Lemma}
\newtheorem{corollary}[theorem]{Corollary}
\newtheorem{definition}[theorem]{Definition}
\newtheorem{remark}[theorem]{Remark}

\title{Task 2: The Weight-Space Nakayama Algebra}
\date{}

\begin{document}
\maketitle

\section{Goal}

This file proves the basic algebraic structure of
\[
A_L^{\mathrm{wt}}=kC_{L+1}/I_L^{\mathrm{wt}}.
\]

The three outputs needed later are:

\begin{itemize}
\item every path of length $>L$ vanishes;
\item the unique surviving length-$L$ path is $p_0=a_L\cdots a_1$;
\item the algebra is a cyclic monomial Nakayama algebra with Kupisch series
  \[
  (L+1,L,\dots,L).
  \]
\end{itemize}

\section{Surviving and vanishing paths}

\begin{lemma}[Length-$L$ paths]
In the quotient algebra $A_L^{\mathrm{wt}}$:

\begin{itemize}
\item the path
  \[
  p_0=a_L\cdots a_1
  \]
  survives;
\item every other oriented path of length $L$ is zero.
\end{itemize}
\end{lemma}

\begin{proof}
The generators of the ideal $I_L^{\mathrm{wt}}$ are exactly the $L$ paths
\[
p_1,\dots,p_L,
\]
which are the length-$L$ oriented paths that pass through the closing arrow
$a_0$.

The remaining length-$L$ path is
\[
p_0=a_L\cdots a_1,
\]
which does not pass through $a_0$.
Since $I_L^{\mathrm{wt}}$ is a monomial ideal generated by the $p_\ell$ for
$1\le \ell\le L$, the class of $p_0$ is not killed in the quotient.
\end{proof}

\begin{lemma}[All paths of length $>L$ vanish]
Every oriented path in $Q=C_{L+1}$ of length at least $L+1$ lies in the ideal
$I_L^{\mathrm{wt}}$.
In particular, every path of length $>L$ is zero in $A_L^{\mathrm{wt}}$.
\end{lemma}

\begin{proof}
Any oriented path of length $L+1$ goes once around the entire cycle. Such a path
contains a contiguous subpath of length $L$ that passes through the closing
arrow $a_0$, hence contains one of the generators $p_1,\dots,p_L$.
Therefore every path of length $L+1$ belongs to $I_L^{\mathrm{wt}}$.

If a path has length $>L+1$, it contains a subpath of length $L+1$, hence also
lies in $I_L^{\mathrm{wt}}$ because $I_L^{\mathrm{wt}}$ is a two-sided ideal.
\end{proof}

\begin{corollary}[Finite-dimensionality]
The algebra $A_L^{\mathrm{wt}}$ is finite-dimensional.
\end{corollary}

\begin{proof}
Only paths of length at most $L$ survive in the quotient, and there are only
finitely many such paths.
\end{proof}

\section{Indecomposable projectives and their lengths}

\begin{definition}[Projective modules]
For each vertex $i\in\{0,\dots,L\}$, let
\[
P(i):=A_L^{\mathrm{wt}}e_i
\]
denote the indecomposable projective corresponding to the vertex idempotent
$e_i$.
\end{definition}

\begin{proposition}[Maximal nonzero paths starting at each vertex]
The maximal nonzero oriented path starting at vertex $i$ has length
\[
\begin{cases}
L, & i=0,\\
L-1, & 1\le i\le L.
\end{cases}
\]
\end{proposition}

\begin{proof}
Starting at $0$, the unique length-$L$ path is
\[
a_L\cdots a_1=p_0,
\]
which survives by the first lemma.

Starting at a vertex $i\in\{1,\dots,L\}$, the unique length-$L$ path from $i$
must pass through the closing arrow $a_0$, hence is exactly the generator
$p_i$. Therefore it vanishes, while every shorter path survives.
\end{proof}

\begin{corollary}[Projective lengths]
The composition lengths of the indecomposable projectives are
\[
\ell(P(0))=L+1,
\qquad
\ell(P(i))=L
\quad
(1\le i\le L).
\]
Equivalently, the Kupisch series is
\[
(c_0,\dots,c_L)=(L+1,L,\dots,L).
\]
\end{corollary}

\begin{remark}
This is the precise difference from the fully truncated cycle algebra used in
\texttt{paper\_v0.tex}, where all length-$L$ paths were killed and the Kupisch series
was uniform.
\end{remark}

\section{Nakayama structure}

\begin{proposition}[Admissibility of the ideal]
For $L\ge 2$, the ideal $I_L^{\mathrm{wt}}$ is admissible:
\[
(kQ)_{\ge L+1}\subseteq I_L^{\mathrm{wt}}\subseteq (kQ)_{\ge 2}.
\]
\end{proposition}

\begin{proof}
The inclusion
\[
I_L^{\mathrm{wt}}\subseteq (kQ)_{\ge 2}
\]
holds because each generator $p_\ell$ has length $L\ge 2$.
The other inclusion is exactly the previous lemma.
\end{proof}

\begin{theorem}[The weight-space algebra is Nakayama]
The algebra $A_L^{\mathrm{wt}}$ is a cyclic monomial Nakayama algebra with
Kupisch series
\[
(L+1,L,\dots,L).
\]
\end{theorem}

\begin{proof}
The ideal $I_L^{\mathrm{wt}}$ is generated by paths, so the quotient is
monomial.

The quiver $Q=C_{L+1}$ is connected and every vertex has exactly one incoming
arrow and one outgoing arrow. Together with admissibility, the standard
bound-quiver criterion implies that the quotient algebra is a cyclic Nakayama
algebra.

The Kupisch series was computed above.
\end{proof}

\section{Deliverables}

This task is complete once the file records:

\begin{itemize}
\item the unique surviving length-$L$ path $p_0=a_L\cdots a_1$;
\item the vanishing of all paths of length $>L$;
\item finite-dimensionality of $A_L^{\mathrm{wt}}$;
\item the projective lengths
  \[
  (L+1,L,\dots,L);
  \]
\item the conclusion that $A_L^{\mathrm{wt}}$ is a cyclic monomial Nakayama
  algebra.
\end{itemize}

\end{document}
